\documentclass[11pt]{article}
\usepackage[utf8]{inputenc}
\usepackage[T1]{fontenc}
\usepackage{lmodern}
\usepackage[margin=1in]{geometry}
\usepackage{amsmath,amssymb}
\usepackage{textcomp}
\usepackage[dvipsnames]{xcolor}
\usepackage{enumitem}
\usepackage{parskip}
\usepackage{titlesec}
\usepackage[colorlinks=true,linkcolor=blue,citecolor=blue,urlcolor=blue]{hyperref}
\usepackage{natbib}

\titleformat{\section}{\large\bfseries}{\thesection}{1em}{}
\titleformat{\subsection}{\normalsize\bfseries}{\thesubsection}{1em}{}
\titleformat{\subsubsection}{\normalsize\bfseries\itshape}{}{0em}{}

\title{\textbf{Contamination Is a Curve, Not a Threshold:\\ Measuring Synthetic-Corpus Dose-Response and Its Detector-Induced Bias}}
\author{}
\date{}

\newcommand{\ph}{\textit{[placeholder --- fill in after experiments are run]}}

\begin{document}
\maketitle
\vspace{-2em}

\begin{quote}
\small\textit{Assembled from the six confirmed stage outputs of an \texttt{auto-research-discussion} run (seed topic B6.2), post research-taste audit (\textcolor{Dandelion}{\textbf{WATCH}} --- one condition fired on the dominant-verb check; fixed by a title/verb reframe), post stage review (2 \textbf{REVISE} findings applied). Status: \textbf{specification / paper draft --- no experiments have been run.} All numbers below are pre-registered design targets, not results.}
\end{quote}

\begin{abstract}
Language models are increasingly pretrained on web-scraped corpora that contain text written by prior-generation language models, but how the \emph{fraction} of such synthetic text maps onto downstream degradation once contamination is partial rather than total remains unmeasured under controlled, confound-isolated conditions. We show that the property making synthetic text detectable --- its excess proximity to the generator's own output distribution --- is the same property through which it erodes the tail mass a trained model needs, so measuring degradation and measuring detectability are two readings of one underlying quantity, and error in one propagates, in a predictable direction, into the other. Starting from a fixed pretrained backbone, we continue pretraining at a fixed token budget across a log-spaced grid of exactly-known synthetic fractions drawn from a provenance-controlled, distribution-matched human/synthetic corpus pair, and score every checkpoint on a pre-registered capability, diversity, and tail-preservation battery. A model-comparison-and-equivalence-testing procedure adjudicates whether degradation follows a flat, linear, convex, or genuine change-point response shape, and a human-text entropy-matched control isolates syntheticness from generic distributional narrowing. We then re-fit the same adjudication pipeline on fractions estimated by a black-box AI-text detector instead of the known ground truth, to measure how much detector error shifts the estimated safety threshold --- and in which direction. We report which response shape dominates for current small-to-mid-scale language models, whether the resulting verdict is stable across capability, diversity, and tail metrics, and what this implies for treating detector output as a corpus-curation policy input.
\end{abstract}

\section{Introduction}

Language models are trained on web-scraped corpora that increasingly contain text those same models, or their predecessors, wrote. A search engine's crawl, a forum thread, a product review --- any of these may now be a machine's paraphrase of an original rather than a human's account of one, and there is no reliable way for a corpus curator to know the exact fraction without running a detector over the data first. The response to this has largely been a warning rather than a measurement: recursive training on generated data eventually collapses a model's output distribution, so contamination is treated as something to minimize, full stop.

Yet the shape of that relationship --- how much degradation per unit of synthetic fraction, measured under controlled conditions --- has not been established. The collapse warning describes what happens at the extreme --- training substantially or entirely on generated data across many generations --- while a theoretical line does already treat the fraction as continuous: scaling-law analyses of collapse predict that even a vanishing synthetic share can break neural scaling in some regimes \citep{TaleOfTails,StrongModelCollapse}, a statistical treatment bounds the maximal tolerable share \citep{SeddikSyntheticStat}, and a separate line derives an optimal human/synthetic mixing ratio --- but all of these live inside simplified regression or scaling-law abstractions, or validate against loss alone, with no confound-isolated empirical dose grid, no capability/diversity/tail battery, and no procedure that can positively certify a safe region rather than assume or deny one. Between these two poles sits the question every curator with a partially contaminated, single-generation corpus actually faces: is there a tolerable fraction below which nothing measurable changes, or does harm accrue from the very first synthetic token, just too slowly to have been reported yet? And if a curator's only practical way to estimate that fraction is to run an off-the-shelf AI-text detector over the corpus, does the resulting estimate tell them where they actually stand, or does the detector's own error quietly move the goalposts?

The nearest existing account of a \emph{closed-form} dose-response-shaped mixing law comes from a different regime entirely. A recent analytical treatment of recursive training on real/synthetic label mixtures derives, in closed form, the mixing weight that minimizes long-run prediction risk and provably prevents collapse --- the closest existing result to a ``dose-response of mixing ratio'' in the literature --- but the analysis lives inside an overparameterized linear or ridge-regression setting, with no language model, no training run, and no empirical dose grid to validate or refute it against. The canonical collapse demonstrations, meanwhile, operate at the opposite pole: they show that indiscriminate or majority-synthetic training degrades a model's distributional tails across recursive generations, and a related result shows that \emph{accumulating} real and synthetic data rather than replacing real data with it avoids this fate --- but both contrast only two discrete regimes, never a continuous fraction, so neither can in principle distinguish a genuine step threshold from gradual, monotonic decline. A third cluster of work treats synthetic data as an asset rather than a threat --- expanding a small corpus into a larger, relationally diverse one for knowledge acquisition, or reweighting synthetic fine-tuning data to better match a real target distribution --- and while informative about when synthetic data helps, none of it varies a synthetic \emph{fraction} systematically or tests for a harm threshold at all. And detection research, which does operationalize ``how much synthetic text is present'' at the level of an individual document, stops exactly there: no detector paper aggregates its scores into a corpus-level contamination estimate or couples that estimate to any downstream measure of training harm.

Turning this gap into a measurable question surfaces four distinct difficulties. Any observed degradation is only attributable to synthetic fraction if every other axis --- token budget, hyperparameters, starting checkpoint, and, critically, the synthetic pool's own document-length and topic distribution --- is held fixed or explicitly matched; otherwise an effect blamed on ``syntheticness'' could just as easily be a generic consequence of feeding the model text that is narrower or differently distributed than usual. Once that confound is controlled, the metrics used to detect degradation must themselves be trustworthy at small model scale: a discontinuous few-shot accuracy score, evaluated on a coarse dose grid, can manufacture an apparent change-point that a continuous metric on the same checkpoints would not show, so the verdict risks being a property of the ruler rather than the model. Even with clean metrics, ``a small but nonzero effect at low dose'' and ``a true zero effect below a threshold'' are statistically adjacent but practically opposite claims, and distinguishing them requires more structure than eyeballing a curve --- an explicit model comparison across candidate response shapes together with a test that can positively certify equivalence to zero, not merely fail to reject it. And because no real-world corpus arrives with its synthetic fraction pre-labeled, any dose-response law derived under known, injected ratios is only actionable in practice if it survives being read through the noisy, document-level detector a curator would actually have to use --- and that detector's error need not simply add noise; it can shift an estimated safety threshold in a specific direction, since the same excess proximity to the generator's distribution that erodes a model's usable tail mass is also what a document-level detector is tuned to notice, most clearly, exactly where the safety question matters least.

Our approach turns each of these four difficulties into a specific, checkable design choice rather than a hedge. A provenance-guaranteed pre-2022 human corpus is paired with a synthetic pool generated to match its document-length and topic distribution, and mixed at document granularity into a log-spaced grid of exactly-known fractions trained from an identical starting checkpoint under an identical budget --- with a human-text control, entropy-matched to the synthetic pool, run alongside it to check whether the resulting curve is attributable to syntheticness specifically or to distributional narrowing in general. Every discontinuous capability metric is paired with a continuous sibling on the same checkpoints, so an apparent threshold can be traced to the model or dismissed as an artifact of metric resolution. A pre-registered model-comparison-and-equivalence-testing procedure adjudicates among flat, linear, convex, and genuine change-point response shapes, replacing a plotted-and-eyeballed verdict with one that must win a stated statistical contest. And the same adjudication pipeline is re-run on fractions estimated by two independent black-box detectors instead of the known ground truth, turning the detector's error --- and its direction --- into a measured quantity rather than an assumed one.

We evaluate this design across a staged pilot-then-scale-up protocol, continuing pretraining of small, pre-ChatGPT-era backbones on the constructed corpus mixtures, and comparing the resulting dose-response fits against a no-history-contamination control, a human-entropy-matched confound control, an accumulate-regime arm, and a single-episode-style detector-calibration check, before testing whether the fitted response shape transfers across backbone scale and synthetic-text generator.

In summary, we contribute:
\begin{itemize}[leftmargin=1.4em]
\item An identification and formalization of the shared-cause asymmetry between contamination detectability and contamination harm --- the same distributional proximity that makes synthetic text recognizable to a detector is what depletes the tail mass a trained model needs --- and a signed prediction (not merely an existence claim) for how detector estimation error biases an estimated safety threshold. The prediction is scoped to the uncorrected, plug-in detector estimate a curator would compute today; a measurement-error-corrected estimator arm is pre-committed to test how much of the bias survives standard errors-in-variables adjustment, isolating the shared-cause component from classical, correctable attenuation.
\item A provenance-controlled, distribution-matched dose-response measurement protocol, including a human-entropy-matched control that isolates syntheticness from generic distributional narrowing as the source of any observed degradation.
\item A model-comparison-and-equivalence-testing adjudication criterion that distinguishes a genuine change-point safety threshold from gradual, resolution-limited monotonic decline, together with an explicit falsification design for each claim it produces.
\item An empirical account of which response shape --- flat, linear, convex, or change-point --- dominates for current small-to-mid-scale language models, whether the verdict is stable across capability, diversity, and tail metrics, and how far a real detector's estimation error shifts it.
\end{itemize}

\section{Problem Definition}

\noindent\textbf{Formal setup.} Fix a pretrained backbone checkpoint $\theta_0$. A mixture corpus $D(\rho)$ of $T$ tokens is constructed for each dose $\rho$ in a log-spaced grid $G \subset [0,1]$, where $\rho$ is the exactly-known fraction of $D(\rho)$ drawn from a synthetic pool rather than a provenance-guaranteed human pool. Continuing pretraining of $\theta_0$ on $D(\rho)$ under seed $s$ yields a checkpoint $\theta(\rho,s)$. A pre-registered metric battery $m \in M$, with one designated primary endpoint $m^*$, is evaluated on every checkpoint; the standardized degradation at dose $\rho$ is
\[
\Delta m(\rho) \;=\; \frac{\operatorname{mean}_s m(\theta(\rho,s)) - \operatorname{mean}_s m(\theta(0,s))}{SD_0(m)},
\]
where $SD_0(m)$ is the seed-level standard deviation of $m$ at $\rho = 0$, signed so that $\Delta m(\rho) > 0$ denotes worse performance. The object of study is the shape of $\Delta m^*(\rho)$ --- whether it is statistically indistinguishable from flat, well-described by a linear or convex power-law trend, or better described by a genuine change-point --- and whether this verdict is stable when $\rho$ is replaced by a detector-estimated $\hat\rho(\rho)$ rather than known exactly.

\noindent\textbf{Why this is hard.} Fixing $\theta_0$, $T$, and every training hyperparameter across the dose grid does not by itself guarantee that $\rho$ is the only systematically varying axis between conditions: the synthetic pool must also be verified, not merely assumed, to match the human pool's document-length and topic distribution, or an observed change in $\Delta m(\rho)$ could reflect generic distributional narrowing rather than synthetic origin specifically --- a confound that requires its own control to rule out, not just a matching step to hope away. Once the data-generating process is trustworthy, ``small but nonzero'' and ``true zero'' are statistically adjacent but decision-relevant opposite hypotheses about $\Delta m^*(\rho)$ at low $\rho$, and adjudicating between a smooth monotone family and a genuine change-point family requires an explicit model-comparison and equivalence-testing procedure, not a plotted curve read by eye. The battery $M$ itself introduces a further hazard: a discontinuous metric evaluated on a coarse dose grid can manufacture an apparent change-point that a continuous metric on the identical checkpoints would not show, so any threshold claim must be checked against a continuous counterpart before it is trusted. Finally, because $\rho$ is exactly known only by experimental construction and never in a real corpus, translating any fitted $\Delta m^*(\rho)$ into an actionable curation policy requires knowing how much a real detector's estimate $\hat\rho(\rho)$ departs from $\rho$, and in which direction --- a question the detection literature, which stops at per-document classification, does not answer.

\section{Method Overview}

At a conceptual level, this is a single measurement pipeline with five cooperating parts. A provenance-controlled corpus pair and dose-grid mixer produce the raw material --- human and synthetic pools distribution-matched to each other, mixed at every tested fraction --- while a fixed-budget continued-pretraining harness turns each mixture into a checkpoint under conditions identical except for the mixture itself. A paired capability/diversity/tail metric battery turns each checkpoint into a set of comparable readings, guarding every discontinuous score with a continuous sibling. A curve-adjudication procedure turns the resulting per-dose readings into a single, statistically defended verdict about the response shape, rather than a plotted impression. And a detector-coupling module re-runs that same adjudication procedure on detector-estimated fractions instead of known ones, turning the practical question --- ``what does this mean for someone who can't inject a known ratio'' --- into a measured, signed quantity.

The \textbf{matched-corpus dose-grid harness} builds a human/synthetic corpus pair matched in document-length and topic distribution, mixes it at document granularity into a log-spaced grid of exactly-known fractions, trains every mixture from an identical starting checkpoint and budget, and runs a human-text entropy-matched control alongside it to isolate syntheticness from generic distributional narrowing. The \textbf{paired metric battery} scores every checkpoint on capability, diversity, and tail/entropy endpoints, pairing every discontinuous score with a continuous counterpart. The \textbf{curve adjudication criterion} takes the resulting per-dose readings and selects among flat, linear, convex, and change-point response families via model comparison and equivalence testing, with a pre-committed interpretation for whichever family wins. The \textbf{detector-coupling module} re-estimates the dose axis with two independent black-box detectors and re-runs the adjudication criterion on the estimated fractions, reporting the resulting shift in the certified-safe region and fitted slope. The \textbf{factorial transfer check} repeats a coarser version of the whole pipeline across backbone scales and synthetic-text generators to test whether the winning response family, not just its coefficients, is stable.

\section{Methodology}

\subsection{Matched-Corpus Dose-Grid Harness}
\textbf{Challenge.} Any observed change in a downstream metric across the dose grid is only attributable to the synthetic fraction if every other axis --- token budget, hyperparameters, starting checkpoint, data-shuffling policy, and the synthetic pool's own distributional properties --- is held fixed or independently verified to match; without this, a ``synthetic contamination'' effect could just as easily be a generic consequence of distributional narrowing.

\textbf{Insight.} Matching the synthetic pool to the human pool and freezing every other training axis turns confound control from an assumption into a checkable gate, but matching alone cannot rule out that any distributionally narrow text --- not specifically synthetic text --- would produce the same curve; only a control that varies distributional entropy independent of syntheticness can isolate the two.

\textbf{Method.} A human reference pool is drawn from a fixed pre-2022 corpus slice, guaranteeing zero synthetic content by provenance. A synthetic pool is generated by prompting an instruction-tuned LLM to regenerate or continue documents drawn from a disjoint slice of the same corpus, with rejection filtering until the pool's document-length and topic-cluster histograms match the human pool's within a pre-registered divergence tolerance. The two pools are mixed at document granularity into a log-spaced grid of exactly-known synthetic fractions and used to continue pretraining an identical starting checkpoint under an identical token budget, hyperparameter set, and shuffling-seed policy for every mixture, with multiple independent seeds per dose. Alongside this grid, a human-text control is constructed by oversampling human documents until the control pool's measured entropy matches the synthetic pool's, and is run through the identical harness at matched doses.

\subsection{Paired Metric Battery}
\textbf{Challenge.} A discontinuous capability metric, evaluated on a coarse dose grid at small model scale where few-shot accuracies are noisy and near-floor, can manufacture an apparent change-point that reflects the metric's own resolution rather than any real property of the model.

\textbf{Insight.} Pairing every discontinuous endpoint with a continuous counterpart evaluated on the identical checkpoints is what lets a later change-point finding be attributed to the model or dismissed as a metric artifact, rather than asserted from whichever metric happened to be reported.

\textbf{Method.} Every checkpoint is scored on held-out clean-corpus perplexity (the sole pre-registered primary endpoint, continuous by construction), few-shot accuracy on a standard small-scale benchmark suite (discontinuous, paired with the perplexity endpoint), generation-based diversity measures (distinct-$n$ statistics, self-similarity, and embedding coverage over sampled continuations), and tail-preservation measures (rare-token recall, tail-decile perplexity, and nearest-neighbor differential entropy of generated text, reusing an indicator established for image-diffusion collapse). Onset dose is defined per metric as the first dose at which a Holm-corrected test detects a significant deviation from the $\rho = 0$ arm.

\subsection{Curve Adjudication Criterion}
\textbf{Challenge.} ``A small but nonzero effect at low dose'' and ``a true zero effect below a threshold'' are statistically adjacent but practically opposite claims about the dose-response curve, and neither a plotted trend nor an unadjusted significance test can reliably tell them apart.

\textbf{Insight.} The only way to distinguish these two hypotheses is to require a candidate change-point family to win an explicit model-comparison contest against smooth alternatives, and to require any ``harmless'' claim to be positively certified by an equivalence test rather than merely surviving a failure to reject the null.

\textbf{Method.} Four candidate response families --- flat, linear, convex power-law, and segmented change-point --- are fit to the seed-level per-metric readings via mixed-effects regression and compared by a corrected information criterion. Change-point existence is additionally tested directly; monotone trend is tested separately; and a dose is certified ``statistically harmless'' only if it passes a two-one-sided equivalence test against the $\rho = 0$ arm under a margin fixed before the confirmatory grid is run, with multiplicity correction applied across the full metric battery. Because seed-level replication per dose is small by design, the equivalence certificate is computed from the pooled mixed-effects fit --- a TOST-style bound on the model-predicted degradation over the candidate safe interval, which borrows strength across the full dose grid --- rather than from per-dose two-sample contrasts, which are underpowered at any margin tighter than several seed-level standard deviations when only a few seeds per dose exist. The equivalence margin itself is finalized only after the pilot stage estimates seed-level variance, and any certificate whose pilot-implied power falls below the pre-registered floor is reported as INCONCLUSIVE rather than converted into evidence of safety. The criterion's interpretation is pre-committed for every possible winner: a convex family win with a non-significant change-point test supports a soft, resolution-limited safe region; a change-point family win with a significant test supports a genuine mechanistic threshold instead, reported as such; and a flat family win even at the highest tested dose supports no measurable harm at the tested scale and routes to a pre-committed reframing of the whole study as an equivalence-certified negative result rather than a silent null finding.

\subsection{Detector-Coupling Module}
\textbf{Challenge.} No real-world corpus arrives with a known synthetic fraction, so any dose-response law fit under exact, experimentally injected ratios is only actionable if it is known how much a real detector's fraction estimate departs from the truth --- and detector error need not behave like simple additive noise on the outcome, because a detector calibrated to separate individual documents is not the same instrument as one calibrated to estimate an aggregate corpus rate.

\textbf{Insight.} The same distributional property that makes synthetic text degrade a trained model --- its excess proximity to the generator's own learned output distribution, which depletes the tail mass a model needs --- is also what a document-level detector is tuned to notice; because that signal is clearest exactly where a mixture is heavily synthetic, a detector's marginal sensitivity is weakest precisely at the low doses where a safety-threshold question is decided, predicting a directionally permissive bias rather than merely a noisy one.

\textbf{Method.} Two independent black-box AI-text-detector instruments are run over a fixed document sample from every mixture in the dose grid, and their per-document scores are aggregated into an estimated corpus-level fraction via a calibration curve fit on the two pure ($\rho = 0$ and $\rho = 1$) pools. Aggregating per-document scores into a corpus-level fraction is a prevalence-estimation (quantification) problem, and naive score averaging is known in that literature to be biased away from extreme prevalences; the module therefore reports both the naive plug-in estimate a practitioner would compute today and a misclassification-corrected variant in the adjusted-classify-and-count style. The signed permissive-bias prediction applies to the naive estimate: classical measurement-error machinery (regression calibration and related errors-in-variables adjustments) can in principle remove attenuation caused by symmetric noise, so whatever bias survives the corrected arm is precisely the shared-cause, dose-dependent component this module exists to isolate. The entire curve-adjudication procedure above is then re-run substituting each detector-estimated fraction for the known one, and the resulting shift in fitted slope, certified-safe region, and winning response family is reported against the ground-truth fit.

\subsection{Factorial Transfer Check}
\textbf{Challenge.} The whole protocol is a small-model, fixed-backbone continued-pretraining proxy; a dose-response law measured at one scale and one synthetic-text generator is a case study, not a policy input, unless there is evidence the same functional family holds elsewhere.

\textbf{Insight.} Curve-shape transfer and coefficient transfer are different claims with different consequences: a curve that keeps its shape but changes steepness across scale is still a usable law with a scale adjustment, whereas a curve whose winning family itself flips means no single law transfers and per-regime calibration is unavoidable.

\textbf{Method.} A coarser dose sub-grid is run as a factorial design crossing two backbone scales and two independently trained synthetic-text generators, and a pooled mixed-effects model estimates the interaction of dose with scale and with generator identity, while the winning response family from the curve-adjudication criterion is checked for exact agreement across all four backbone-by-generator cells, independent of whether the fitted coefficients themselves agree.

\section{Experiments}

\subsection{Experimental Setup}

\textbf{Datasets.} A pre-2022 slice of C4 or The Pile serves as the human reference pool, contamination-free by provenance and consistent with the pretraining-era provenance of small pre-ChatGPT backbones such as the Pythia suite. Synthetic pools are generated by an instruction-tuned LLM (a second, independently trained generator is added at the scale-up stage specifically to test whether the fitted curve family is a property of the phenomenon or an artifact of one generator's output distribution). WikiText-103 is used as a second, independently curated clean held-out corpus alongside the in-domain human validation slice, so that the primary-endpoint perplexity trend is checked against a source whose vocabulary and register differ from the training pool's own held-out split --- guarding against a curve-shape verdict that is an artifact of overfitting to one validation source's idiosyncrasies. A standard small-scale few-shot benchmark suite provides the discontinuous capability endpoints paired against the continuous perplexity primary endpoint.

\textbf{Baselines} (sanity-check and confound-control comparators, not competitors to beat --- this is a measurement study, and no research question here is framed as ``our procedure beats baseline X''). A $\rho = 0\%$ clean continued-pretraining arm is the necessary control for every dose-level claim, separating dose effects from continued-pretraining drift. A no-continued-pretraining frozen-backbone anchor isolates how much continued pretraining itself, independent of any synthetic content, moves each metric. A human-text entropy-matched arm is the rival-explanation control: human text oversampled until its measured entropy matches the synthetic pool's, run at the same doses, to test whether the measured effect is attributable to syntheticness specifically or to distributional narrowing generally. An accumulate-regime arm, in which synthetic text is added on top of the full human set rather than replacing part of it, connects the classic replace-versus-accumulate contrast to this design's confound-controlled continuous dose axis (prior empirical work contrasts accumulate/replace workflows at discrete synthetic quantities without provenance-matched confound control; \citealp{KazdanCollapseThrive}). An entropy-maximizing data-selection arm at a fixed intermediate dose serves as a secondary mitigation reference.

\textbf{Metrics.} The standardized degradation $\Delta m(\rho)$ on every battery endpoint; per-metric onset dose; AICc weights across the four candidate response families; the Davies change-point test statistic; the Jonckheere--Terpstra monotone-trend statistic; the TOST-certified safe region; the detector calibration error and the resulting slope-attenuation ratio and safe-region shift; and cross-cell agreement of the winning response family in the factorial transfer check.

\textbf{Research questions:}
\begin{enumerate}[leftmargin=1.6em]
\item RQ1 --- Does the measured degradation actually trace to synthetic origin, or could a generically distribution-narrowed human control produce the same curve?
\item RQ2 --- Do capability, diversity, and tail/entropy endpoints agree on when degradation begins, or does an apparent change-point on a discontinuous metric dissolve under its paired continuous counterpart?
\item RQ3 --- Which response family --- flat, linear, convex, or genuine change-point --- best describes the primary-endpoint dose-response curve, and is any resulting ``safe region'' resolution-limited or mechanistic?
\item RQ4 --- When the dose axis is estimated by a real detector instead of known exactly, how far and in which direction does the certified-safe region and fitted slope shift?
\item RQ5 --- Is the winning response family stable across backbone scale and synthetic-text generator?
\end{enumerate}

\subsubsection{RQ1: Does the measured degradation trace to synthetic origin, or to generic distributional narrowing?}
\textbf{Setup.} At matched doses, the primary-endpoint dose-response curve fit on the synthetic-fraction arm is compared against the curve fit on the human-text entropy-matched control arm, via an interaction term in the pooled mixed-effects model (endpoint $\sim$ dose $\times$ arm-type).

\textbf{Expected outcome.} The synthetic-fraction and entropy-matched-human curves are statistically distinguishable (significant dose $\times$ arm-type interaction), with the synthetic arm degrading more steeply at matched entropy --- demonstrating that the measured effect is attributable to syntheticness specifically, not to distributional narrowing of any kind.

\textbf{Falsification condition.} Refuted if the two curves are statistically indistinguishable (non-significant interaction, overlapping fits) --- in which case the finding is not void, but must be re-scoped from ``synthetic text harms'' to the weaker, still-informative claim that ``entropy-depleted text harms, and synthetic text is one instance of it,'' exactly the pre-committed re-scoping this control exists to trigger rather than paper over.

\textbf{Hypothesis (plain language).} The degradation this paper measures is caused by synthetic origin specifically, not merely by the kind of distributional narrowing any low-diversity text could produce.

\textbf{Metrics.} Dose $\times$ arm-type interaction coefficient and its confidence interval on the primary endpoint; curve-fit overlap between the two arms.

\textbf{Result.} \ph

\subsubsection{RQ2: Do capability, diversity, and tail metrics agree on when degradation begins?}
\textbf{Setup.} Per-metric onset dose (the first Holm-corrected significant deviation from $\rho = 0$) is computed for every endpoint in the battery and compared across capability, diversity, and tail/entropy families; any discontinuous metric with a significant change-point test is checked against its paired continuous counterpart on the identical checkpoints.

\textbf{Expected outcome.} Diversity and tail/entropy endpoints show a significantly earlier onset dose than capability endpoints, and any change-point detected on a discontinuous few-shot-accuracy metric fails to replicate on its paired continuous metric --- indicating that a threshold read off only a capability benchmark would understate how early real degradation begins.

\textbf{Falsification condition.} Refuted if all metric families yield statistically indistinguishable onset doses and curve shapes, or if capability endpoints degrade before diversity/tail endpoints --- either result would mean a single cheap endpoint (clean perplexity) already captures the full dose-response picture, which is itself a useful, reportable simplification rather than a failed hypothesis.

\textbf{Hypothesis (plain language).} A curator who monitors only benchmark accuracy would certify a corpus as safe at a contamination level that has already measurably narrowed what the model can say.

\textbf{Metrics.} Onset dose per metric family; paired discontinuous/continuous change-point agreement.

\textbf{Result.} \ph

\subsubsection{RQ3: Which response family describes the dose-response curve, and is any safe region resolution-limited or mechanistic?}
\textbf{Setup.} On the primary endpoint, across every tested backbone $\times$ domain condition cell, the four candidate response families are fit via mixed-effects regression and compared by AICc; the Davies change-point test and Jonckheere--Terpstra monotone-trend test are applied; a model-based TOST equivalence test against $\rho = 0$ (computed on the pooled mixed-effects fit over the candidate safe interval, with a pilot-calibrated pre-registered margin; see Methodology) certifies the safe region.

\textbf{Expected outcome.} The convex power-law family wins the AICc comparison with a fitted exponent significantly greater than one, the change-point test is non-significant, and the certified-safe region is non-empty but confined to the lowest one or two grid doses --- a soft, resolution-limited safe region rather than a genuine mechanistic threshold.

\textbf{Falsification condition.} Refuted toward a genuine threshold if the segmented change-point family instead wins the AICc comparison with a significant Davies test --- in which case the finding is not weaker but a stronger and different actionable result (a real safety ceiling at the fitted change-point location), reported as such rather than forced into the convex-family framing. Refuted toward no measurable harm if the flat family wins even at the highest tested dose, in which case the study is reframed as a pre-committed equivalence-certified negative result rather than a silent null.

\textbf{Hypothesis (plain language).} Degradation is real and accelerating with dose, but the low-dose ``safe'' region reflects the resolution of what can currently be measured, not a protective mechanism.

\textbf{Metrics.} AICc weights across the four families; fitted exponent and its confidence interval; Davies test $p$-value; TOST-certified safe-region boundary.

\textbf{Result.} \ph

\subsubsection{RQ4: How far, and in which direction, does detector estimation error shift the certified-safe region?}
\textbf{Setup.} Two independent black-box detector instruments score a fixed document sample from every mixture; scores are mapped to an estimated corpus fraction via a calibration curve fit on the two pure pools; the full curve-adjudication procedure from RQ3 is re-run substituting the estimated fraction for the known one.

\textbf{Expected outcome.} Re-fitting on the detector-estimated fraction attenuates the fitted primary-endpoint slope and widens the certified-safe region relative to the ground-truth fit, with the calibration error negative (under-estimating the true fraction) at low doses --- detector noise biases a contamination policy toward permissiveness, not merely toward imprecision.

\textbf{Falsification condition.} Refuted if the detector-estimated fit does not meaningfully attenuate the slope or shift the safe region (estimated fraction tracks the truth closely enough that thresholds move by less than one grid step) --- a positive, practically useful result meaning off-the-shelf detectors can be used directly as a policy input today. Also refuted in direction if the shift is instead toward over-alarming (a narrower certified-safe region) rather than permissiveness --- in which case the signed direction, not the existence of a shift, is what the paper reports as overturned.

\textbf{Hypothesis (plain language).} Treating a detector's fraction estimate as a safety-policy input without correction makes contamination look safer than it actually is, not just noisier.

\textbf{Metrics.} Calibration error (estimated minus true fraction) as a function of true fraction; slope-attenuation ratio between detector-estimated and ground-truth fits; safe-region shift in grid steps.

\textbf{Result.} \ph

\subsubsection{RQ5: Is the winning response family stable across backbone scale and synthetic-text generator?}
\textbf{Setup.} A coarser dose sub-grid is run as a factorial design crossing two backbone scales and two independently trained synthetic-text generators; a pooled mixed-effects model estimates dose $\times$ scale and dose $\times$ generator interaction terms, and the AICc-winning response family is compared for exact agreement across all four cells.

\textbf{Expected outcome.} The winning response family on the primary endpoint is identical across all four backbone-scale $\times$ generator cells, while the fitted dose-coefficient magnitude differs significantly by scale --- the curve's shape transfers even though its steepness does not, supporting a scale-adjustment rather than a per-regime law.

\textbf{Falsification condition.} Refuted if the winning response family itself differs across any of the four cells (e.g., a change-point win at one scale but a convex win at another) --- in which case no single dose-response law transfers, contamination tolerances must be calibrated per training regime rather than once, and the released measurement harness becomes the deliverable for that per-regime calibration rather than a single portable curve.

\textbf{Hypothesis (plain language).} The \emph{shape} of the danger --- accelerating, threshold, or flat --- is a property of the phenomenon that travels across scale and generator, even if the exact steepness does not.

\textbf{Metrics.} Per-cell AICc-winning family (categorical agreement); dose $\times$ scale and dose $\times$ generator interaction coefficients from the pooled model.

\textbf{Result.} \ph

\vspace{0.5em}
\noindent\textbf{Further evaluation.} With more time, a natural extension is a small from-scratch pretraining replication at a reduced scale, to directly check whether the continued-pretraining proxy this design relies on understates or overstates degradation relative to contamination present from the first pretraining token --- explicitly out of scope here because it multiplies the compute budget by the number of dose points rather than reusing one fixed starting checkpoint, and would need its own confound-control design (isolating budget from regime) rather than a direct extension of the matched-corpus harness above. A second extension, absorbed as a design note from a sibling study on domain-selective contamination effects: stratifying the dose grid by domain and token rarity --- recent evidence that per-parameter memorization concentrates on rare tokens rather than scaling with corpus size suggests contamination harm may likewise concentrate where human tail mass is thinnest, making a domain/rarity-stratified dose-response a follow-on axis of this harness rather than a separate study.

\section{Related Work}

The closest prior work to this paper is an analytical mixing-ratio result for recursive training on real/synthetic label mixtures \citep{PreventCollapseMixing}, and it is worth stating the delta plainly before surveying the wider landscape: that work already derives, in closed form, the real/synthetic mixing weight that minimizes long-run prediction risk and provably prevents collapse --- the closest existing ``dose-response of mixing ratio'' result in the literature --- but the analysis lives entirely inside an overparameterized linear or ridge-regression abstraction, with no language model trained, no empirical dose grid swept, and no downstream capability or diversity endpoint measured. This paper's contribution is precisely to take that theoretical shape and ask whether, and in what form, it holds when the ``model'' is an actual language model and the ``mixing ratio'' is a variable a real corpus curator has to estimate rather than set exogenously.

A large-scale empirical study closes part of that theory-to-practice gap from the opposite direction \citep{DemystifySynthPretrain}: training over 1,000 language models under a unified protocol, it sweeps synthetic-fraction mixture ratios directly and reports an empirically ``good'' ratio ($\sim$30\% for rephrased synthetic data), while explicitly rejecting the generality of prior model-collapse conjectures for one-round (n=1) training on rephrased synthetic data. This is the closest empirical precedent to this paper's core manipulated variable, but it is benefit/efficiency-framed --- optimizing for validation-loss convergence speed via scaling-law curve-fitting, not a pre-registered statistical test for a harm threshold --- and its own stated limitations explicitly decline multi-generational dynamics, fine-grained capability/diversity batteries, and any formal change-point test, which is precisely the machinery this paper adds on top of the same manipulated variable.

\textbf{Mixing-ratio and data-mixture theory.} Beyond the closest work above, a companion analysis \citep{CuratedSelfConsuming} models curated retraining as implicit preference optimization, explaining both alignment gains and diversity collapse under repeated curation --- but curation, not a synthetic-fraction dose axis, is the manipulated variable there. Separately, pretraining data-mixture optimization work \citep{DoReMi} establishes that domain-weight proportions causally and transferably affect language-model quality across scales, using a small proxy model and a distributionally robust objective to set weights that transfer to a much larger training run; a very recent causal-inference reformulation of the same domain-mixture problem \citep{CAUSALMIX1104} targets the same transfer-without-new-proxy-runs goal. All three lines establish that mixture proportions matter and can be optimized, but none treats a synthetic-versus-human axis as the object of study, and none reports a degradation endpoint at all --- the mixture-optimization machinery targets speed of convergence to a good domain mix, not the location of a harm threshold.

\textbf{Model collapse at the extreme dose.} \citet{ShumailovCurseRecursion} show that indiscriminate training on model-generated data progressively erodes a distribution's tails across recursive generations, \citet{GerstgrasserAccumulate} show this can be avoided if synthetic data accumulates alongside real data rather than replacing it, and a diffusion-model study \citep{CloserLookCollapse} reframes collapse as an entropy-driven generalization-to-memorization transition and proposes entropy-maximizing selection as a countermeasure. A follow-up study co-authored by one of the field's foundational authors \citep{IterativeDeploymentImpro4940} complicates the picture further, showing curated iterative self-deployment can \emph{improve} planning skills across generations while noting ``model collapse may still occur, or curation may only delay it'' --- a rival, curation-contingent narrative this paper's accumulate-regime control arm is designed to engage with. All of these are foundational to how this paper conceives of ``harm,'' and this paper's tail-preservation and entropy endpoints are directly adapted from the diffusion-collapse indicator. But all operate at the extreme or at a binary regime choice: the manipulated variable is generation count or a binary replace-versus-accumulate/curate choice, never a continuous within-generation fraction, so none can distinguish a genuine step threshold from gradual monotonic decline --- precisely the adjudication this paper's curve-comparison criterion is built to make.

\textbf{Fraction-continuous collapse analyses.} A further line --- closest of all to this paper's dose \emph{axis}, and conceded here plainly --- treats the synthetic proportion as a continuous variable inside a scaling-law or statistical abstraction: model collapse recast as a change of scaling laws \citep{TaleOfTails}, the ``strong collapse'' result that even very small synthetic fractions can degrade downstream scaling behavior \citep{StrongModelCollapse}, a statistical bound on the maximal tolerable synthetic share \citep{SeddikSyntheticStat}, and an empirical study of accumulate-versus-replace workflows at varying synthetic quantities \citep{KazdanCollapseThrive}. These works establish that the continuous-fraction question is well-posed and give it \emph{rival} theoretical answers --- from ``no safe fraction exists'' to ``a bounded tolerable share exists.'' What none of them provides is what this design supplies: an exogenously injected, provenance-guaranteed dose grid on an actual language model with an entropy-matched human confound control, a pre-registered adjudication able to decide \emph{between} those rival answers (a strong-collapse world is, in this design's terms, a linear- or convex-family win with an empty certified-safe region --- a pre-committed rival hypothesis, not an assumption), and any coupling of the dose axis to the detector a practitioner would actually use to estimate it.

\textbf{Synthetic data as an asset, and controlled-intervention infrastructure.} A separate, benefit-framed line of work \citep{SyntheticCPT} expands small domain corpora into larger, relationally diverse synthetic corpora for knowledge acquisition, and a related reweighting scheme \citep{NotAllLLMData} corrects for distribution mismatch between synthetic and real data in classification fine-tuning. Both demonstrate that synthetic data can help under the right conditions, but neither varies a synthetic fraction systematically or tests for a harm threshold, and the reweighting work in particular is scoped to fine-tuning and classification rather than pretraining. Infrastructure for cheap controlled pretraining interventions \citep{TrainOnceAnswerAll} --- embedding many small, independently probed data manipulations into a single shared training run --- is the methodological precedent this paper's harness borrows its controlled-intervention discipline from, but that infrastructure's own design premise restricts each intervention to a small fraction of total tokens, structurally excluding the high-dose mixture sweeps this paper's grid requires.

\textbf{AI-text detection.} Black-box detection methods, including a recoverability-style approach \citep{DandR} that shuffles and re-orders text within an LLM call and scores how well the model restores the original order, operationalize ``how much synthetic text is present'' at the level of an individual document. This paper reuses such a detector directly as one of its two coupling instruments, but the detection literature itself stops at per-document classification: no detector paper aggregates its scores into a corpus-level contamination estimate, or couples that estimate to any measure of downstream training harm --- the gap this paper's detector-coupling module is specifically built to close.

\section{Conclusion}

This work reframes corpus contamination from a binary warning into a measurable function: contamination has a shape, not just a ceiling, and the instrument used to estimate it in practice is itself biased in a specific, predictable direction. The central methodological wager is that a dose-response curve is only meaningful once syntheticness is isolated from generic distributional narrowing, every discontinuous metric is checked against a continuous sibling, and a change-point claim is required to win an explicit statistical contest rather than being read off a plot --- weaker designs risk mistaking metric artifacts or narrowing effects for the specific phenomenon this paper measures. Whichever response shape dominates for current language models, the finding is decision-relevant either way: a soft, resolution-limited safe region means today's low-dose contamination is not yet visibly dangerous but is not zero either, while a genuine change-point would hand curators an actual, defensible ceiling. Two limitations bound this claim from the outset and are natural next steps rather than afterthoughts: the entire protocol is a small-model, fixed-backbone continued-pretraining proxy for production-scale pretraining-from-scratch, and while the factorial transfer check reports whether the fitted response family holds across the tested scales and generators, it cannot by construction certify transfer beyond them --- a risk sharpened by recent evidence that standard small-proxy-model protocols can mis-rank data recipes relative to fully-tuned large-scale outcomes \citep{CanSmallRunsGuideCuration}, which this paper's factorial transfer check is designed to detect rather than assume away; and the detector-coupling analysis characterizes estimation bias under the two specific instruments tested here, not a fundamental limit on detection --- a materially better detector could narrow, though the shared-cause argument in this paper's Introduction suggests it may not eliminate, the permissive bias this paper measures.

\vspace{1em}
\hrule
\vspace{1em}

\section*{Pre-handoff checklist}
\begin{itemize}[leftmargin=1.4em]
\item \textbf{Organization/navigation} --- narrative flows problem $\to$ method $\to$ experiments $\to$ related work without forward references to undefined terms.
\item \textbf{Terminology discipline} --- ``dose-response curve,'' ``certified-safe region,'' ``response family,'' ``onset dose,'' ``calibration error'' each defined at first use in Problem Definition/Methodology, used consistently thereafter; no internal stage IDs (\texttt{sp\_*}, \texttt{m\_*}, \texttt{rq\_*}) leaked into prose --- checked directly against the discussion pipeline log.
\item \textbf{Notation with worked example} --- the formal setup section defines $(\theta_0, \rho, G, D(\rho), \theta(\rho,s), M, m^*, \Delta m(\rho), SD_0(m))$ and the Problem Definition/Methodology sections consistently reuse this notation; no separate worked numeric example is included since no numbers exist yet (would require fabricating results).
\item \textbf{RQ framing matches outcomes} --- all 5 RQs are graded/quantitative (slopes, onset doses, AICc weights, calibration error), each answered in its own subsection with an explicit falsifiable expected outcome and a distinct falsification condition, not a binary yes/no.
\item \textbf{No typesetting artifacts} --- only explicit placeholder markers in Results, per the drafting convention this is a pre-registered spec, not a report.
\item \textbf{Artifact plan} --- the matched-corpus generation pipeline, training configs, checkpoint evaluation tables, curve-fitting code, and detector-calibration data are the natural release artifacts; not yet built, so not claimed as available.
\end{itemize}

\section*{Hand-off}
This is a draft, not a finished paper --- experiments have not run. Natural next steps: prose polish and venue-specific LaTeX formatting once a target venue is chosen, or revising an earlier stage if piloting surfaces a problem (e.g., if the pilot stage surfaces that seed-level noise makes RQ3's low-dose resolution far coarser than expected).

\bibliographystyle{plainnat}
\bibliography{refs}

\end{document}
